\chapter{The renderer and the user interface}
\label{ch:renderer}

\section{Shaders that can be reloaded}
\code{ReloadShaders} compiles the four entry points from their files with
\code{D3DCompileFromFile}, collecting the compiler's messages in a log. It creates the new
shader objects and input layouts first, and replaces the old ones only if everything succeeded.
A typing mistake in a shader therefore leaves the program running with the previous version.
The input layouts use \code{offsetof} on the C++ vertex structures, so the structures are the
only description of the vertex format.

\excerpt{reloadshaders}

\section{The off-screen target}
The 3D view is rendered into a texture the size of the Viewport panel, which Dear ImGui then
shows like any other image. The target uses four samples per pixel (multisample antialiasing),
which smooths the silhouette of the solid and the thin lines of the probe. It is written through
an sRGB view, so the GPU encodes the linear colours of the lighting, and \emph{resolved} into a
single-sample texture that the interface reads through a plain view. This is the same
typeless-texture technique used for the colour maps (Section~\ref{sec:typeless}).

\excerpt{scenetarget}

\section{Drawing the scene}
\code{RenderScene} clears the target, uploads the constant buffer (with \code{WRITE\_DISCARD}, so
the CPU never waits for the GPU to finish the previous frame), and draws the surface with one
indexed draw call. The four maps are bound to both shader stages: the vertex shader needs the
displacement map. A missing map is replaced by a neutral grey texture.

The lines of the probe and of the light marker are drawn twice from the same buffer: first at
30\,\% opacity without the depth test, so that parts hidden behind the solid remain faintly
visible, then normally with the depth test. Finally the textures are unbound, because Direct3D
refuses to bind a resource for reading and writing at once, and the multisampled image is
resolved. The resolve uses the sRGB format, so the samples are averaged in linear light.

\excerpt{renderscene}

\section{The frame of the application}
\code{App::Frame} is the function of Figure~\ref{fig:frame}. The order matters in one respect:
the probe is decided while the map panels are built, before the viewport, which needs the
probe's result to draw the arrows.

\excerpt{frame}

\section{Images that answer the mouse}
Every image of a map in the interface (the tiles of the Maps panel and the large image of the
Inspector) is drawn by \code{App::MapImage}, which also converts the mouse position into texture
coordinates:
\[
  (u,v) = \Bigl(\frac{p_x-m_x}{W},\ \frac{p_y-m_y}{H}\Bigr)\ \text{clamped to }[0,1]^2,
\]
where $\vect m$ is the image's top-left corner on the screen and $(W,H)$ its size. Hovering moves
the probe, a left click pins it, a right click releases it. The probe's crosshair is drawn on
every map image, so the eye can compare the maps at the same texel.

\excerpt{mapimageui}

\section{Filling the constant buffer}
The Viewport panel computes the camera matrices (right-handed, looking at the origin) and copies
the application's state into the constant buffer. A map is used only if it is switched on
\emph{and} present, and the mipmap level for the displacement is chosen to match the mesh
resolution:

\excerpt{fillconstants}

\section{The main loop}
\file{main.cpp} creates the window, initializes Dear ImGui (with docking, and with its layout file
next to the executable), and runs the classic loop: process the window messages, measure the
elapsed time, build the interface (which renders the scene off-screen), draw the interface into
the window, and present it synchronized with the display.

\excerpt{mainloop}
